% Generated deterministically from the audited Markdown source.
% Responsible author: Carptopus.
\documentclass[11pt]{article}
\usepackage[a4paper,margin=27mm]{geometry}
\usepackage[T1]{fontenc}
\usepackage[utf8]{inputenc}
\usepackage{lmodern}
\usepackage{microtype}
\usepackage{amsmath,amssymb,amsthm,mathtools}
\usepackage{needspace}
\usepackage{xcolor}
\usepackage[colorlinks=true,linkcolor=blue!55!black,citecolor=blue!55!black,urlcolor=blue!65!black]{hyperref}
\usepackage{xurl}
\usepackage{fancyhdr}

\pagestyle{fancy}
\fancyhf{}
\fancyhead[L]{Carptopus}
\fancyhead[R]{Cographic surface operations}
\fancyfoot[C]{\thepage}
\setlength{\headheight}{14pt}
\setlength{\parindent}{0pt}
\setlength{\parskip}{0.44em}
\setlength{\emergencystretch}{4em}
\urlstyle{same}

\title{Triangle sums, peripheral cycles, and contractions\\
for orderable cographic matroids}
\author{Carptopus\\\href{mailto:carptopus@163.com}{carptopus@163.com}}
\date{4 September 2026}

\hypersetup{
  pdftitle={Triangle sums, peripheral cycles, and contractions for orderable cographic matroids},
  pdfauthor={Carptopus},
  pdfsubject={Structural operations for face-adjacency orderings of cographic matroids},
  pdfkeywords={orderable matroid, cographic matroid, surface triangulation, triangle sum, edge contraction, peripheral cycle}
}

\begin{document}
\maketitle

\begin{center}
\fcolorbox{black!35}{black!4}{\parbox{0.88\textwidth}{\centering\small
Internally verified candidate manuscript, version 0.1-beta. The mathematical proof and
complete manuscript have passed project-internal zero-trust audits. External mathematical
review and formal peer review are pending.}}
\end{center}

\begin{abstract}


For a finite simplicial triangulation $T$ of a closed surface, let
$\Lambda_T$ be the graph on the edges of the $1$-skeleton in which two edges
are adjacent when they lie in a common triangular face. We study the property
$U(T)$ that $\Lambda_T[B]$ is connected for every bond $B$ of the
$1$-skeleton. The induced graph is automatically $2$-regular, so $U(T)$
provides a consistent ordering of the circuits of the cographic matroid.

We first prove a graph-theoretic gluing theorem for cographic matroids with a
compatible rooted triangle. In the surface setting this becomes an exact
connected-sum theorem:

\[
U(T_1\mathbin{\#_\triangle}T_2)
\quad\Longleftrightarrow\quad
U(T_1)\ \text{and}\ U(T_2).
\]

Together with the projective-plane theorem proved in an earlier paper, this
gives infinitely many $3$-connected, regular, non-graphic, orderable
cographic matroids on every closed nonorientable surface. We then establish
two complementary structural restrictions. A nonfacial induced cycle whose
vertex deletion leaves the graph connected must be one-sided, and $U(T)$
descends under every legal edge contraction. Consequently, existence on a
fixed surface reduces to its finite set of irreducible triangulations. We also
give an explicit Klein-bottle example showing that $U$ is not preserved by
general vertex splitting. These results provide composition, obstruction, and
reduction mechanisms without claiming a classification of all orderable
cographic matroids or all surface triangulations satisfying $U$.

\end{abstract}

\section{Introduction}


Crenshaw and Oxley introduced \emph{orderable matroids}: matroids whose circuits
admit reversible cyclic orderings that agree on every shared adjacent pair
\cite{crenshaw-oxley-published}. Graphic matroids are naturally orderable, while the binary case is
structurally difficult. Their current author manuscript proves that every
$4$-connected regular orderable matroid is graphic \cite[Theorem~1.5]{crenshaw-oxley-v12}. The exact
$3$-connected boundary is different. An earlier paper \cite{carptopus-rp2} proved that every
finite simplicial triangulation of the real projective plane gives an
orderable cographic matroid, yielding infinitely many $3$-connected regular
non-graphic examples.

The earlier proof used a particular global adjacency relation: two graph
edges are adjacent when they meet in a triangular face. Restricted to any
bond, this relation is a disjoint union of cycles, and a topological argument
on the projective plane forces a single cycle. The present paper asks which
operations preserve this stronger, specified ordering and what its failure
reveals about the topology of a triangulation.

The results have three layers. First, Section 3 gives a graph-theoretic
triangle-rooted gluing theorem that is not restricted to surface
triangulations. Second, Section 4 proves that the surface property is preserved
and reflected by simplicial connected sum along a triangular face. This
extends \cite{carptopus-rp2} from the projective plane to explicit infinite families on every
closed nonorientable surface. Third, Sections 5 and 6 give a topological
obstruction and a contraction theorem. The obstruction limits peripheral
cycles in any positive example; the contraction theorem reduces the existence
question on each fixed surface to finitely many irreducible triangulations.

The specified property studied here is stronger than abstract orderability:
$U(T)$ implies that $M^*(T^{(1)})$ is orderable, but an abstract consistent
ordering need not arise from the triangular faces of the given embedding. We
keep this distinction throughout.

\section{Definitions and the face-adjacency property}


All graphs and simplicial complexes are finite. If $G$ is connected, the
circuits of its cographic matroid $M^*(G)$ are the bonds of $G$. Equivalently,
for a nonempty proper set $S\subset V(G)$, the cut $\delta_G(S)$ is a bond
exactly when both $G[S]$ and $G[V(G)\setminus S]$ are connected.

We use the following sufficient form of consistency.

\textbf{Lemma 2.1 (global adjacency).} Let $M$ be a matroid. Suppose that a simple
graph $A$ on $E(M)$ has the property that $A[C]$ is a cycle for every circuit
$C$ of $M$. Reading each cycle in either direction gives a consistent family
of reversible cyclic orderings. Hence $M$ is orderable.

\textbf{Proof.} Whether two elements are adjacent is determined by the single
global graph $A$. Thus, whenever two circuits contain the same pair, that pair
is adjacent in one induced ordering if and only if it is adjacent in the
other. $\square$

Let $T$ be a simplicial triangulation of a closed connected surface, let
$G=T^{(1)}$, and define the \emph{face-adjacency graph} $\Lambda_T$ on $E(G)$ by
joining two distinct graph edges when they belong to a common triangular
face. For a cut $B$, each triangular face contains either zero or two edges
of $B$. Moreover, each cut edge lies in two faces. It follows that
$\Lambda_T[B]$ is $2$-regular.

We say that $T$ has the \emph{single-boundary property}, denoted by $U(T)$, when
$\Lambda_T[B]$ is connected for every bond $B$ of $G$. Thus $U(T)$ is
equivalent to the statement that this particular face-adjacency relation
induces a cycle on every bond. By Lemma 2.1,

\[
U(T)\quad\Longrightarrow\quad M^*(T^{(1)})\text{ is orderable}.
\]

There is also a useful topological interpretation. Give the vertices on the
two sides of a cut the values $0$ and $1$, extend affinely over each simplex,
and take the level set at $1/2$. Its segments through the mixed faces form a
disjoint union of circles, and their combinatorial adjacency graph is exactly
$\Lambda_T[B]$. Thus $U(T)$ says that the level set associated with every
bond has one component.

\section{A triangle-rooted gluing theorem}


The first result is independent of surface topology.

Let $H$ be a graph containing a triangle $R$ with edges $r_1,r_2,r_3$.
Suppose $M^*(H)$ has a consistent ordering represented by a global adjacency
graph $A_H$ as in Lemma 2.1. We call this ordering *compatible with the rooted
triangle $R$* when

\[
A_H[\{r_1,r_2,r_3\}]=K_3.
\]

This does not require $R$ to be a bond or a face in a specified embedding.

\textbf{Theorem 3.1 (triangle-rooted gluing).} Let $H_1$ and $H_2$ be finite simple
$3$-connected graphs whose intersection is exactly a common triangle $R$:

\[
V(H_1)\cap V(H_2)=V(R),
\qquad
E(H_1)\cap E(H_2)=E(R).
\]

Assume that each $M^*(H_i)$ has a global consistent adjacency graph $A_i$
compatible with $R$. If $G=H_1\cup_R H_2$, with the three edges of $R$
retained, then $M^*(G)$ is orderable.

\textbf{Proof.} Form a graph $A$ on $E(G)$ by taking $A_1\cup A_2$ and deleting
the three adjacencies between pairs of root edges. We prove that $A[B]$ is a
cycle for every bond $B$ of $G$.

Use the two sides of $B$ to color $V(G)$ with two colors. If the three
vertices of $R$ are monochromatic, the opposite color cannot occur in the
interiors of both factors, since those vertices would have no same-colored
route between the factors. In the active factor, every component of the
root-colored side meets $R$, and the monochromatic triangle $R$ joins all
such components. Thus $B$ lies entirely in one factor, where it is a bond and
contains no root edge. The deletion of root-root adjacencies does not affect
it, so $A[B]=A_i[B]$ is a cycle.

Suppose instead that $R$ is bichromatic. Exactly two root edges, say $r_a$
and $r_b$, cross the cut. In either factor, every component induced by one
color must meet $R$, since otherwise it could not connect through the other
factor. The same-colored vertices of $R$ are adjacent within the triangle, so
each color induces a connected subgraph in each factor. Therefore
$B_i=B\cap E(H_i)$ is a bond of $H_i$ and contains $r_a,r_b$.

Compatibility makes $r_ar_b$ an edge of the cycle $A_i[B_i]$. Deleting that
adjacency leaves a Hamilton path $P_i$ from $r_a$ to $r_b$. Since a
$3$-connected graph has no bond of size two, each path contains a non-root
element. The paths have only their endpoints in common, and
$P_1\cup P_2=A[B]$ is a cycle through every element of $B$. Lemma 2.1 now
proves the result. $\square$

If $H_i$ is the $1$-skeleton of a closed surface triangulation satisfying
$U$, the face-adjacency graph supplies such a rooted ordering at every
triangular face. The theorem also applies beyond triangulations. For example,
a $3$-connected plane graph with a triangular face has a compatible cographic
ordering coming from its planar dual and can be glued to any other compatible
factor.

\section{Simplicial triangle sums}


Let $T_1,T_2$ be closed connected simplicial triangulations, with chosen
triangular faces $F_1,F_2$. Delete the interiors of those faces and identify
their boundary triangles by a vertex bijection. Assume that the interiors of
the two complexes are otherwise disjoint. Write the resulting triangulation
as $T_1\mathbin{\#_\triangle}T_2$. The identified triangle is retained in the
$1$-skeleton but is no longer a face; we call it the \emph{seam}.

\textbf{Theorem 4.1 (triangle-sum equivalence).}

\[
U(T_1\mathbin{\#_\triangle}T_2)
\quad\Longleftrightarrow\quad
U(T_1)\ \text{and}\ U(T_2).
\]

\textbf{Proof.} The forward construction in Theorem 3.1 applies to the
face-adjacency graphs of the factors. Equivalently, consider a bond in the
sum. If the seam vertices are monochromatic, the bond lies in one factor. If
they are bichromatic, its restriction to either factor is a bond containing
the same two seam edges. In each capped factor the corresponding boundary is
one circle. Removing the chosen face turns that circle into an arc between
the midpoints of the two crossing seam edges, and the two arcs glue to one
circle. Hence $U(T_1)$ and $U(T_2)$ imply $U(T_1\mathbin{\#_\triangle}T_2)$.

For the converse, suppose that $T_1$ fails $U$. Choose a bond of its
$1$-skeleton whose face-adjacency boundary has $k\geq2$ components and use
its sides to color the vertices. If the chosen face is monochromatic, give
all non-seam vertices of $T_2$ that same color. The cut remains a bond in the
sum and retains all $k$ boundary components.

If the chosen face is bichromatic, give every non-seam vertex of $T_2$ the
majority color on the seam. The minority-colored side in $T_2$ is a single
seam vertex. The other side is connected because the $1$-skeleton of a closed
simplicial triangulation is $3$-connected and remains connected after that
vertex is deleted. Thus the extended cut is again a bond. In the punctured
$T_2$, the boundary around the minority seam vertex is one arc. It glues only
to the component of the original boundary that crossed the deleted face; the
remaining $k-1$ circles persist. The sum therefore fails $U$. The same
argument applies with the factors reversed. $\square$

The equivalence iterates along any tree of triangle sums. In particular, a
triangulation decomposed along separating nonfacial triangles satisfies $U$
if and only if every capped leaf factor does. The statement does not require
uniqueness of the chosen decomposition tree.

\textbf{Corollary 4.2 (triangular disk replacement).} Let $D$ be a simplicial
triangulated disk with triangular boundary, whose interior vertices and edges
are disjoint from a closed triangulation $T$. Replacing a triangular face of
$T$ by $D$ preserves and reflects $U$.

\textbf{Proof.} If $D$ is one triangle there is nothing to prove. Otherwise, cap
$D$ with a new triangular face. The result is a sphere triangulation, which
satisfies $U$: in the planar dual, every bond of the primal graph corresponds
to one dual cycle, with exactly the face-adjacency order. Apply Theorem 4.1.
$\square$

\subsection{Infinite families in every nonorientable genus}


Let $N_h$ denote the closed nonorientable surface with $h$ crosscaps. The
earlier projective-plane theorem \cite{carptopus-rp2} says that every finite simplicial
triangulation of $N_1=\mathbb{RP}^2$ satisfies $U$.

\textbf{Theorem 4.3.} For every integer $h\geq1$, take any $h$ finite simplicial
triangulations of $\mathbb{RP}^2$ and join them successively along triangular
faces. The resulting triangulation $T$ of $N_h$ satisfies $U(T)$. For every
$h$, this construction yields infinitely many pairwise nonisomorphic
$3$-connected, regular, non-graphic, orderable cographic matroids.

\textbf{Proof.} The first assertion follows by repeated application of Theorem
4.1. The $1$-skeleton of a closed simplicial triangulation is $3$-connected.
It can also be seen directly here: after deleting at most two vertices, each
factor remains connected and every seam retains a vertex joining its two
sides.

If the triangulation has $n$ vertices and $m$ edges, Euler's formula gives

\[
m=3n-6+3h>3n-6.
\]

Its $1$-skeleton is therefore nonplanar. Its cographic matroid is regular and
binary, is $3$-connected, and is non-graphic by the graphic-cographic
planarity criterion for $3$-connected graphs. It is orderable by $U(T)$ and
Lemma 2.1. Finally, repeated stellar subdivision of a triangular face is a
special case of Corollary 4.2 and increases the ground-set size, giving
infinitely many pairwise nonisomorphic matroids. $\square$

\section{A peripheral-cycle obstruction and a failed inverse operation}


The connected-sum theorem supplies positive constructions. The next result is
a necessary condition for an arbitrary triangulation to satisfy $U$.

\textbf{Theorem 5.1 (peripheral-cycle obstruction).} Let $T$ be a closed connected
simplicial triangulation satisfying $U(T)$, with $G=T^{(1)}$. If $C$ is an
induced cycle of $G$, $G-V(C)$ is connected, and $C$ is not a triangular face
of $T$, then $C$ is one-sided in the surface.

\textbf{Proof.} Since $G[V(C)]=C$ and $G-V(C)$ are connected,
$\delta_G(V(C))$ is a bond. Because $C$ is induced and nonfacial, the full
subcomplex of $T$ induced by $V(C)$ is the circle $C$, rather than a disk.
In every simplex, the lower half of the two-valued piecewise-linear function
associated with the cut collapses linearly onto its face spanned by the
vertices of $C$. These simplexwise collapses agree on common faces. The
standard derived-neighborhood criterion therefore identifies the lower half
as a PL regular neighborhood of the full induced subcomplex $C$ \cite{mohar-thomassen}. Its
boundary components are precisely the components of
$\Lambda_T[\delta_G(V(C))]$.

If $C$ is two-sided, its regular neighborhood is an annulus and has two
boundary components, contradicting $U(T)$. A one-sided cycle has a Möbius-band
neighborhood and one boundary component, so it is the only remaining
possibility. $\square$

\textbf{Corollary 5.2.} Every two-sided nonfacial induced cycle in a triangulation
satisfying $U$ is vertex-separating. In particular, every two-sided nonfacial
triangle is a $3$-vertex separator.

The inverse of edge contraction is vertex splitting. It is natural to ask
whether arbitrary legal vertex splits preserve $U$, which would turn finite
irreducible triangulations into generators for all positive examples. The
answer is no.

\textbf{Proposition 5.3.} The property $U$ is not preserved by general legal vertex
splitting, even for $3$-connected triangulations of the Klein bottle.

\textbf{Proof.} Begin with the following nine-vertex Klein-bottle triangulation,
obtained as the triangle sum of two six-vertex projective-plane
triangulations:

\begin{verbatim}
014 017 024 025 035 038 067 068 123
125 136 145 168 178 234 345 367 378
\end{verbatim}


It satisfies $U$ by Theorem 4.1. The link of vertex $0$ in cyclic order is

\begin{verbatim}
1 4 2 5 3 8 6 7.
\end{verbatim}


Split the two link arcs at positions $1$ and $5$, introducing vertex $9$.
The resulting faces are

\begin{verbatim}
014 017 049 067 068 089 123 125 136 145
168 178 234 249 259 345 359 367 378 389.
\end{verbatim}


This is a closed $3$-connected Klein-bottle triangulation and has no
separating nonfacial triangle. For $S=\{0,3,4,6\}$, the cut
$B=\delta(S)$ is a bond. Its $17$ edges are partitioned by face adjacency
into the following two cycles:

\begin{verbatim}
01 07 67 37 38 39 35 45 14
08 09 49 24 23 13 16 68
\end{verbatim}


Thus $\Lambda_T[B]$ is disconnected and the split triangulation fails $U$.
$\square$

An exact exhaustive check of all $90$ nontrivial single splits of the same
labelled base triangulation gives $63$ positive operations retaining a
separating nonfacial triangle and $27$ negative operations with no such
triangle. These operations produce $78$ distinct labelled face sets: $52$
positive and $26$ negative. This finite table is a calibration around one
base example, not a classification up to isomorphism and not a general
equivalence between seams and $U$.

\section{Descent under legal edge contraction}


Although $U$ need not survive the inverse operation, it always descends under
legal edge contraction.

Let $e=uv$ be an edge of a closed simplicial triangulation $T$ satisfying the
usual link condition, so that contracting $e$ gives another simplicial
triangulation $T'=T/e$ of the same surface. Let $w$ be the image of $u,v$.

\textbf{Theorem 6.1 (contraction descent).}

\[
U(T)\quad\Longrightarrow\quad U(T/e).
\]

\textbf{Proof.} Let $B'=\delta_{T'^{(1)}}(S')$ be a bond and color its two sides.
Lift the coloring to $T$ by giving both $u$ and $v$ the color of $w$. The
resulting cut $B$ is a bond. Indeed, a monochromatic path in $T'$ lifts
directly except possibly when it enters $w$ through an edge incident only
with $u$ and leaves through an edge incident only with $v$; in that case the
monochromatic edge $uv$ joins the two lifted pieces. Hence both color classes
remain connected.

By the link condition, the common neighbors of $u$ and $v$ are exactly the
third vertices $a,b$ of the two faces incident with $uv$. Away from $uv$, the
only edge identifications under contraction are

\[
ua,va\mapsto wa,
\qquad
ub,vb\mapsto wb.
\]

The full link condition also prevents two distinct surviving faces from
colliding under these identifications. Thus, after the two incident faces
$uva$ and $uvb$ disappear, every other face of $T$ corresponds uniquely to a
face of $T'$.

For $z\in\{a,b\}$ of the same color as $u,v$, neither $uz$ nor $vz$ belongs
to $B$, so no cut-edge identification occurs. If $z$ has the opposite color,
then $uz,vz\in B$ and those two edges are adjacent in $\Lambda_T[B]$ through
the disappearing mixed face $uvz$; contraction replaces them by the single
edge $wz$. All other mixed faces and their induced adjacencies correspond
bijectively before and after contraction. Consequently
$\Lambda_{T'}[B']$ is obtained from the connected graph $\Lambda_T[B]$ by
contracting zero, one, or two of these specified adjacency edges. It is
therefore connected. As in Section 2 it is also $2$-regular, so it is a single
cycle. Since $B'$ was arbitrary, $U(T')$ holds. $\square$

A triangulation is \emph{irreducible} if no edge can be legally contracted while
remaining a triangulation of the same surface.

\textbf{Corollary 6.2 (finite irreducible reduction).} Every closed surface
triangulation satisfying $U$ contracts through $U$-preserving steps to an
irreducible triangulation of the same surface. Hence a fixed closed surface
has a triangulation satisfying $U$ if and only if one of its irreducible
triangulations satisfies $U$.

\textbf{Proof.} Repeatedly apply Theorem 6.1 until no legal contraction remains.
The reverse implication in the existence statement takes the irreducible
triangulation itself. Each fixed closed surface has only finitely many
irreducible triangulations \cite{joret-wood}. $\square$

The corollary is an existence reduction, not a classification of all positive
triangulations. Proposition 5.3 shows exactly why the reduction cannot simply
be reversed by allowing arbitrary vertex splits. Complete irreducible lists
are known for several low-genus surfaces \cite{sulanke}, but those enumerations are not
used in the proof.

\section{Verification and scope}


The source package contains three separate finite checks. They are independent
of the general proofs but share a common exact implementation of graph bonds
and face adjacency:

\begin{enumerate}
\item they build iterated triangle sums of the six-vertex projective-plane
   triangulation and exhaust all bonds through nonorientable genus five;
\item they reconstruct Proposition 5.3, check the topology and
   $3$-connectivity of every split result, and exhaust every bond in all $90$
   split operations;
\item they check a bounded collection of two-step splits as a diagnostic for a
   possible stronger seam theorem.
\end{enumerate}


The corresponding entry points are
\texttt{verify\_nonorientable\_triangle\_sum\_family.py},
\texttt{check\_vertex\_split\_nonclosure.py}, and
\texttt{probe\_two\_step\_seam\_persistence.py}. The first two support statements
recorded in this manuscript. The third returned no positive prime example in
its bounded search, but that \texttt{NO\_HIT} is not used as evidence for an
existence, nonexistence, or classification theorem. The proofs of Theorems
3.1, 4.1, 4.3, 5.1, and 6.1 are general and do not depend on the finite
checks. The release artifact will record the exact commands, shared
dependency, and file hashes.

We make no claim that every orderable cographic matroid arises from a surface
triangulation, that every abstract ordering is a face-adjacency ordering, or
that the prime triangulations satisfying $U$ have been classified. We also do
not claim that the rooted compatibility condition in Theorem 3.1 is
necessary. Every constructed example contains triangular cycles, which give
three-element cocircuits and exact $3$-separations; for $h\geq2$, the seam
triangles give additional such separations. The examples therefore do not
contradict the $4$-connected result of Crenshaw and Oxley \cite{crenshaw-oxley-v12}.

The current literature check covered the published paper \cite{crenshaw-oxley-published}, the latest
author manuscript \cite{crenshaw-oxley-v12}, the preceding projective-plane construction \cite{carptopus-rp2}, and
the standard irreducible-triangulation sources \cite{sulanke,joret-wood}. No direct statement of
the triangle-rooted gluing theorem, the triangle-sum equivalence, the
peripheral obstruction, or the contraction-descent theorem was located. This
records the search boundary; it is not a claim of exhaustive worldwide
priority. The present manuscript is an internally developed candidate and has
not undergone external peer review.

\Needspace{0.16\textheight}
\section*{AI-assisted research disclosure}


Carptopus is the responsible author and takes responsibility for the claims,
proofs, source use, and released artifacts. OpenAI Codex was used for
AI-assisted literature organization, proof development and stress testing,
exact verification, adversarial auditing, and manuscript preparation.

\begin{thebibliography}{9}
\footnotesize
\raggedright

\bibitem{crenshaw-oxley-published}
C.~Crenshaw and J.~Oxley,
\emph{Ordering circuits of matroids},
Electronic Journal of Combinatorics \textbf{29} (2022), no.~4, Paper P4.31.
\href{https://doi.org/10.37236/11117}{doi:10.37236/11117}.

\bibitem{crenshaw-oxley-v12}
C.~Crenshaw and J.~Oxley,
\emph{Ordering circuits of matroids}, author manuscript, version 12.
\url{https://www.math.lsu.edu/~oxley/Orderability_v12.pdf}.

\bibitem{carptopus-rp2}
Carptopus,
\emph{Projective-plane triangulations yield orderable cographic matroids},
version 0.1-beta (2026).
\href{https://doi.org/10.5281/zenodo.22165775}{doi:10.5281/zenodo.22165775}.

\bibitem{sulanke}
T.~Sulanke,
\emph{Irreducible triangulations of low genus surfaces},
arXiv:math/0606690 (2006).
\url{https://arxiv.org/abs/math/0606690}.

\bibitem{joret-wood}
G.~Joret and D.~R.~Wood,
\emph{Irreducible triangulations are small},
arXiv:0907.1421 (2009).
\url{https://arxiv.org/abs/0907.1421}.

\bibitem{mohar-thomassen}
B.~Mohar and C.~Thomassen,
\emph{Graphs on Surfaces}, Johns Hopkins University Press, 2001.

\end{thebibliography}

\end{document}
